\section{Green-Trace Calibration and the Sign-Free Completion Route}
\label{sec:GreenTraceCalibration}

This section records the most economical way to turn the roadmap into a full initial-data proof. The point is to stop trying to upgrade the variational Euler condition to the pointwise sign
\(\tr_\Sigma k\ge0\). The no-go theorem in Appendix~\ref{subsec:GapClosureMath} shows that such a pointwise upgrade cannot follow from a bare scalar adjoint-cone argument. The replacement is to use the adjoint-cone condition at exactly the place where the proof needs it: the ADM mass variation in the conformal correction.

\subsection{The calibration principle}

Let \(\Sigma=\Sigma_{\max}\) be the maximum-area trapped surface supplied by the compactness part of the program. The constrained first variation gives
\[
    L_\Sigma^{-*}(\tr_\Sigma k)\le0
\]
in the distributional sense, with the outward-normal convention fixed in Appendix~\ref{subsec:GapClosureMath}. Equivalently,
\begin{equation}\label{eq:AdjointConePairing37}
    \int_\Sigma (\tr_\Sigma k)\,w_\eta\,dA \le 0
    \qquad\text{for every }\eta\ge0,\quad L_\Sigma w_\eta=\eta .
\end{equation}
The direct favorable-jump proof used the stronger pointwise condition \(\tr_\Sigma k\ge0\) to make the corner scalar-curvature measure nonnegative. The sign-free route asks for less: it asks that the particular test functions which determine the ADM mass correction lie in the cone
\[
    \mathcal C_\Sigma:=\{w_\eta:\eta\ge0,\ L_\Sigma w_\eta=\eta\}.
\]
Then the variational inequality already gives the correct mass sign.

\begin{definition}[Green-trace calibration]\label{def:GreenTraceCalibration}
Let \((\tM,\tg)\) be the sealed Jang manifold associated with \(\Sigma\), and let \(\mathcal L_{\tg}\) denote the conformal Laplacian appearing in the smoothing-correction step,
\[
    \mathcal L_{\tg}u=-8\Delta_{\tg}u+R_{\tg}^{\mathrm{reg}}u
\]
away from the interface, with the interface measure \(2(\tr_\Sigma k)\delta_\Sigma\) kept as a source term. Let \(G_\infty\) be the positive Green potential normalized by
\[
    G_\infty(x)=|x|^{-1}+O(|x|^{-2})
    \qquad\text{on the asymptotically flat end}.
\]
We say that \(\Sigma\) satisfies the \emph{Green-trace calibration} if the trace
\[
    \Gamma_\Sigma:=G_\infty|_\Sigma
\]
belongs to the closed MOTS admissible cone \(\overline{\mathcal C_\Sigma}\). Equivalently, there exist nonnegative \(\eta_j\in C^\infty(\Sigma)\) such that
\[
    L_\Sigma^{-1}\eta_j\to \Gamma_\Sigma
\]
in \(H^1(\Sigma)\).
\end{definition}

\begin{remark}
This condition is deliberately weaker than \(\tr_\Sigma k\ge0\). It is a one-function calibration, or after localization a finite-dimensional family of calibrations, tailored to the mass functional. It is therefore not excluded by Theorem~\ref{thm:NoGoVariational}, which rules out only a blind upgrade from the adjoint-cone law to pointwise positivity of the source.
\end{remark}

\subsection{Mass variation under the calibrated sign-free correction}

\begin{lemma}[Interface contribution to the ADM coefficient]\label{lem:InterfaceADMGreen}
For the conformal correction \(u_\epsilon\to u\) used to remove the negative part of the raw smoothed scalar curvature, the first variation of the ADM mass contributed by the interface measure is
\begin{equation}\label{eq:ADMGreenPairing37}
    \delta m_\Sigma
    =
    c_0\int_\Sigma (\tr_\Sigma k)\,\Gamma_\Sigma\,dA
    + o(1),
\end{equation}
where \(c_0>0\) is the universal normalization constant determined by the Green function convention. The \(o(1)\) term vanishes as the smoothing scale tends to zero.
\end{lemma}

\begin{proof}
Let \(v_\epsilon=u_\epsilon-1\). Linearizing the conformal equation around \(u=1\) gives
\[
    -8\Delta_{\hat g_\epsilon}v_\epsilon
    =
    -R_{\hat g_\epsilon}+ \text{quadratic error}.
\]
The scalar curvature of the raw smoothing converges as a measure to
\[
    R_{\tg}^{\mathrm{reg}}+2(\tr_\Sigma k)\delta_\Sigma,
\]
with \(R_{\hat g_\epsilon}^-\to0\) in \(L^{3/2}\) after the controlled correction estimates of Appendix~\ref{app:InternalSmoothing}. Pairing the equation with the normalized Green potential \(G_\infty\), integrating by parts on large coordinate balls, and then letting the radius tend to infinity identifies the \(r^{-1}\)-coefficient of \(v_\epsilon\), hence the ADM mass variation. The regular bulk term is already controlled by the Bray--Khuri divergence identity and is nonpositive in the mass chain. The singular interface term contributes exactly the pairing in~\eqref{eq:ADMGreenPairing37}. The smoothing error tends to zero by the \(L^{3/2}\)-smallness of the negative part and the uniform \(L^6\) Sobolev bound for \(G_\infty\) on the collar.
\end{proof}

\begin{proposition}[Adjoint-cone compensation]\label{prop:AdjointConeCompensation}
Assume \(\Sigma\) is a constrained maximum-area MOTS satisfying the adjoint-cone Euler condition~\eqref{eq:AdjointConePairing37}. If the Green-trace calibration of Definition~\ref{def:GreenTraceCalibration} holds, then
\[
    \delta m_\Sigma\le0.
\]
Consequently the sign-changing part of \(\tr_\Sigma k\) does not obstruct the conformal mass comparison.
\end{proposition}

\begin{proof}
By Green-trace calibration, choose \(\eta_j\ge0\) with \(w_j=L_\Sigma^{-1}\eta_j\to\Gamma_\Sigma\) in \(H^1(\Sigma)\). Since \(\tr_\Sigma k\) is smooth on the MOTS, the pairings converge:
\[
    \int_\Sigma(\tr_\Sigma k)\Gamma_\Sigma\,dA
    =
    \lim_{j\to\infty}\int_\Sigma(\tr_\Sigma k)w_j\,dA .
\]
Each term on the right is nonpositive by the adjoint-cone Euler condition. Hence the limit is nonpositive. Lemma~\ref{lem:InterfaceADMGreen} then gives \(\delta m_\Sigma\le0\).
\end{proof}

\subsection{Completion theorem}

\begin{theorem}[Completion theorem from Green-trace calibration]\label{thm:FullPenroseFromCalibration}
Let \((M,g,k)\) be asymptotically flat initial data satisfying DEC, and let \(\Sigma_0\) be a closed future trapped surface. Suppose:
\begin{enumerate}
    \item the compactness theorem for maximum-area trapped surfaces holds, producing a smooth MOTS \(\Sigma_{\max}\) with \(A(\Sigma_{\max})\ge A(\Sigma_0)\);
    \item \(\Sigma_{\max}\) satisfies the adjoint-cone Euler condition of Theorem~\ref{thm:AdjointConeEuler};
    \item \(\Sigma_{\max}\) satisfies the Green-trace calibration of Definition~\ref{def:GreenTraceCalibration};
    \item the standard smoothing, conformal correction, Mosco convergence, and AMO limit estimates hold as in Sections~\ref{sec:AMO}--\ref{sec:Synthesis}.
\end{enumerate}
Then
\[
    M_{\ADM}(g)\ge \sqrt{\frac{A(\Sigma_0)}{16\pi}}.
\]
\end{theorem}

\begin{proof}
The compactness theorem supplies a comparison MOTS \(\Sigma_{\max}\) with
\[
    A(\Sigma_{\max})\ge A(\Sigma_0).
\]
Run the Jang deformation and sealing construction with blow-up along \(\Sigma_{\max}\). The only place where the earlier direct proof required \(\tr_{\Sigma_{\max}}k\ge0\) pointwise was the mass control in the corner correction. Proposition~\ref{prop:AdjointConeCompensation} replaces that pointwise sign by the calibrated adjoint-cone pairing and shows that the interface contribution to the ADM mass is nonpositive. The remaining bulk terms are controlled exactly as in the direct MOTS proof: DEC gives the nonnegative square terms in the Jang identity, the Bray--Khuri divergence identity gives the conformal mass comparison, and the conformally corrected smooth approximants satisfy the AMO hypotheses. Therefore
\[
    M_{\ADM}(g)
    \ge M_{\ADM}(\tg)
    \ge \sqrt{\frac{A(\Sigma_{\max})}{16\pi}}
    \ge \sqrt{\frac{A(\Sigma_0)}{16\pi}}.
\]
\end{proof}

\subsection{How to prove the calibration}

The remaining work is now sharply localized. One does not need to prove \(\tr_\Sigma k\ge0\). One needs to prove that the trace of the ADM Green potential belongs to the admissible cone of the MOTS stability operator. This is the concrete next theorem to pursue.

\begin{conjecture}[Green-trace calibration conjecture]\label{conj:GreenTraceCalibration}
Let \(\Sigma_{\max}\) be a smooth strictly stable constrained maximum-area MOTS in asymptotically flat initial data satisfying DEC. Then the trace \(\Gamma_\Sigma=G_\infty|_\Sigma\) of the normalized ADM Green potential satisfies
\[
    L_\Sigma\Gamma_\Sigma\ge0
\]
in the weak sense on \(\Sigma\). Equivalently, \(\Gamma_\Sigma\in\overline{\mathcal C_\Sigma}\).
\end{conjecture}

\begin{remark}[Why this is the right next target]
The conjecture is a boundary-to-horizon positivity statement, not a pointwise sign statement for \(\tr_\Sigma k\). It couples the asymptotic mass functional to the horizon stability operator. This is precisely the geometric information missing from the bare variational no-go theorem.
\end{remark}

\subsection{Boundary identity for the Green trace}

We next rewrite the conjecture as a concrete boundary estimate. Write the MOTS stability operator in the standard form
\begin{equation}\label{eq:MOTSStabilityGeneral37}
    L_\Sigma \psi
    =
    -\Delta_\Sigma\psi+2X\cdot\nabla_\Sigma\psi+Q_\Sigma\psi,
\end{equation}
where \(X\) is the tangential connection one-form induced by \(k\) and \(Q_\Sigma\) is the zeroth-order stability potential. Let \(\nu\) denote the normal to the interface pointing toward the asymptotically flat side in the sealed Jang manifold. Near \(\Sigma\), decompose the regular part of the conformal Laplacian as
\[
    \mathcal L_{\tg}^{\mathrm{reg}}
    =
    -8\left(\partial_\nu^2+H_\Sigma\partial_\nu+\Delta_\Sigma\right)
    +R_{\tg}^{\mathrm{reg}}
    +\mathcal E_{\mathrm{collar}},
\]
where \(\mathcal E_{\mathrm{collar}}\) contains the remaining first-order and lower-order collar terms after the displayed normal/tangential splitting. In what follows its sign is fixed by this displayed decomposition; equivalently, it may be read as the signed residual needed for the boundary identity below.

\begin{definition}[Green flux defect]\label{def:GreenFluxDefect}
For the normalized ADM Green potential \(G_\infty\), define the horizon flux defect
\begin{equation}\label{eq:GreenFluxDefect37}
    \mathfrak D_\Sigma(G_\infty)
    :=
    \partial_\nu^2G_\infty
    +H_\Sigma\partial_\nu G_\infty
    +2X\cdot\nabla_\Sigma\Gamma_\Sigma
    +\left(Q_\Sigma-\frac18 R_{\tg}^{\mathrm{reg}}\big|_\Sigma\right)\Gamma_\Sigma
    -\frac18\mathcal E_{\mathrm{collar}}G_\infty\big|_\Sigma .
\end{equation}
\end{definition}

\begin{lemma}[Boundary identity for the calibration]\label{lem:BoundaryIdentityCalibration}
Assume \(G_\infty\) is \(C^2\) up to the regular side of \(\Sigma\) and solves \(\mathcal L_{\tg}^{\mathrm{reg}}G_\infty=0\) there. Then
\begin{equation}\label{eq:BoundaryIdentityCalibration}
    L_\Sigma\Gamma_\Sigma=\mathfrak D_\Sigma(G_\infty)
\end{equation}
in the weak sense on \(\Sigma\).
\end{lemma}

\begin{proof}
Restrict the equation \(\mathcal L_{\tg}^{\mathrm{reg}}G_\infty=0\) to the collar. The displayed decomposition gives
\[
    -\Delta_\Sigma\Gamma_\Sigma
    =
    \partial_\nu^2G_\infty
    +H_\Sigma\partial_\nu G_\infty
    -\frac18 R_{\tg}^{\mathrm{reg}}\Gamma_\Sigma
    -\frac18\mathcal E_{\mathrm{collar}}G_\infty .
\]
Adding the drift and stability-potential terms \(2X\cdot\nabla_\Sigma\Gamma_\Sigma+Q_\Sigma\Gamma_\Sigma\) gives exactly~\eqref{eq:GreenFluxDefect37}. The weak formulation follows by testing against \(C^\infty(\Sigma)\) functions and using the trace theorem in the collar.
\end{proof}

\begin{corollary}[Flux-defect criterion]\label{cor:FluxDefectCriterion}
If
\[
    \mathfrak D_\Sigma(G_\infty)\ge0
\]
weakly on \(\Sigma\), then the Green-trace calibration holds.
\end{corollary}

\begin{proof}
By Lemma~\ref{lem:BoundaryIdentityCalibration}, \(L_\Sigma\Gamma_\Sigma=\mathfrak D_\Sigma(G_\infty)\ge0\). Strict stability gives the maximum principle and positivity of \(L_\Sigma^{-1}\) on nonnegative data. Hence \(\Gamma_\Sigma=L_\Sigma^{-1}\mathfrak D_\Sigma(G_\infty)\) belongs to \(\overline{\mathcal C_\Sigma}\).
\end{proof}

\begin{proposition}[Model verification on an exact product cylinder]\label{prop:ProductCylinderCalibration}
Assume the sealed Jang end is exactly cylindrical near \(\Sigma\), the tangential drift terms in \(L_\Sigma\) vanish, and the Green potential is constant on the cross-sections of the cylinder. Then the Green-trace calibration holds.
\end{proposition}

\begin{proof}
In the exact product case, \(G_\infty|_\Sigma\) is a positive constant \(c\). The MOTS stability operator reduces to
\[
    L_\Sigma=-\Delta_\Sigma+V_\Sigma
\]
with \(V_\Sigma\ge0\) under the strict stability/product-cylinder hypotheses. Hence
\[
    L_\Sigma(c)=cV_\Sigma\ge0.
\]
Thus \(\Gamma_\Sigma=c=L_\Sigma^{-1}(cV_\Sigma)\) lies in the admissible cone.
\end{proof}

\begin{theorem}[Perturbative Green-trace calibration]\label{thm:PerturbativeGreenTraceCalibration}
Let \(\Sigma\) be strictly stable, and suppose the collar geometry is \(C^{2,\alpha}\)-close to an exact product-cylinder model for which
\[
    L_\Sigma^{(0)}\Gamma_\Sigma^{(0)}\ge \kappa_0
\]
for some \(\kappa_0>0\). More precisely, assume
\[
    \|L_\Sigma-L_\Sigma^{(0)}\|_{C^{1,\alpha}\to C^{0,\alpha}}
    +
    \|\Gamma_\Sigma-\Gamma_\Sigma^{(0)}\|_{C^{2,\alpha}}
    \le \varepsilon .
\]
Then for \(\varepsilon\le\varepsilon_0(\kappa_0,\|L_\Sigma^{(0)}\|)\), the Green-trace calibration holds.
\end{theorem}

\begin{proof}
Compute
\[
    L_\Sigma\Gamma_\Sigma
    =
    L_\Sigma^{(0)}\Gamma_\Sigma^{(0)}
    +(L_\Sigma-L_\Sigma^{(0)})\Gamma_\Sigma^{(0)}
    +L_\Sigma(\Gamma_\Sigma-\Gamma_\Sigma^{(0)}).
\]
The first term is at least \(\kappa_0\). The two error terms are bounded in \(C^{0,\alpha}\) by \(C\varepsilon\), where \(C\) depends only on the product model and the strict-stability ellipticity constants. Choosing \(\varepsilon_0\le\kappa_0/(2C)\) gives
\[
    L_\Sigma\Gamma_\Sigma\ge \kappa_0/2>0.
\]
The flux-defect criterion then places \(\Gamma_\Sigma\) in the admissible cone.
\end{proof}

\subsection{Logarithmic flux and a Riccati barrier}
\label{subsec:LogFluxRiccati}

The flux-defect conjecture can be sharpened further by replacing the second normal derivative of \(G_\infty\) with a first-order logarithmic flux. Since \(G_\infty>0\), set
\[
    \Phi:=\log G_\infty,\qquad
    \Upsilon:=\Phi|_\Sigma=\log\Gamma_\Sigma,\qquad
    \rho:=\partial_\nu\Phi|_\Sigma
    =\frac{\partial_\nu G_\infty}{\Gamma_\Sigma}.
\]
Then
\[
    \partial_\nu G_\infty=\rho\Gamma_\Sigma,\qquad
    \partial_\nu^2G_\infty=(\partial_\nu\rho+\rho^2)\Gamma_\Sigma,
    \qquad
    \nabla_\Sigma\Gamma_\Sigma=\Gamma_\Sigma\nabla_\Sigma\Upsilon .
\]
Substituting these identities into Definition~\ref{def:GreenFluxDefect} gives the exact logarithmic form
\begin{equation}\label{eq:LogFluxDefect37}
    \mathfrak D_\Sigma(G_\infty)
    =
    \Gamma_\Sigma
    \left[
        \partial_\nu\rho+\rho^2+H_\Sigma\rho
        +2X\cdot\nabla_\Sigma\Upsilon
        +Q_\Sigma-\frac18 R_{\tg}^{\mathrm{reg}}\big|_\Sigma
        -\frac18\mathcal E_{\mathrm{collar}}\mathbf 1
    \right],
\end{equation}
where \(\mathcal E_{\mathrm{collar}}\mathbf 1\) denotes the zeroth-order value of the signed collar residual after applying \(\mathcal E_{\mathrm{collar}}\) to \(G_\infty\) and dividing by \(\Gamma_\Sigma\). If \(\mathcal E_{\mathrm{collar}}\) has first-order tangential or normal parts, those terms are included in this shorthand after writing them in logarithmic variables.

The regular conformal equation also gives a Riccati equation for \(\rho\). Dividing
\(\mathcal L_{\tg}^{\mathrm{reg}}G_\infty=0\) by \(G_\infty\) in the collar yields
\begin{equation}\label{eq:GreenRiccati37}
    \partial_\nu\rho+\rho^2+H_\Sigma\rho
    =
    -\Delta_\Sigma\Upsilon
    -|\nabla_\Sigma\Upsilon|^2
    +\frac18R_{\tg}^{\mathrm{reg}}\big|_\Sigma
    +\frac18\mathcal E_{\mathrm{collar}}^{\log}(\Phi),
\end{equation}
with \(\mathcal E_{\mathrm{collar}}^{\log}(\Phi)\) the logarithmic collar residual. This equation is useful in two complementary directions: \eqref{eq:LogFluxDefect37} is a direct lower-bound criterion for the defect, while \eqref{eq:GreenRiccati37} converts the same condition into a tangential inequality for \(\Upsilon\).

\begin{definition}[Riccati barrier inequality]\label{def:RiccatiBarrier37}
We say that the ADM Green potential satisfies the Riccati barrier inequality on \(\Sigma\) if
\begin{equation}\label{eq:RiccatiBarrier37}
    \partial_\nu\rho+\rho^2+H_\Sigma\rho
    +2X\cdot\nabla_\Sigma\Upsilon
    +Q_\Sigma-\frac18 R_{\tg}^{\mathrm{reg}}\big|_\Sigma
    -\frac18\mathcal E_{\mathrm{collar}}\mathbf 1
    \ge0
\end{equation}
weakly on \(\Sigma\).
\end{definition}

\begin{proposition}[Riccati barrier implies calibration]\label{prop:RiccatiBarrierCalibration}
If the Riccati barrier inequality~\eqref{eq:RiccatiBarrier37} holds, then \(\Sigma\) satisfies the Green-trace calibration.
\end{proposition}

\begin{proof}
Because \(G_\infty>0\), the trace \(\Gamma_\Sigma\) is positive. Equation~\eqref{eq:LogFluxDefect37} shows that the flux defect is the product of \(\Gamma_\Sigma\) with the left side of~\eqref{eq:RiccatiBarrier37}. Hence
\[
    \mathfrak D_\Sigma(G_\infty)\ge0 .
\]
The flux-defect criterion, Corollary~\ref{cor:FluxDefectCriterion}, then gives the Green-trace calibration.
\end{proof}

\begin{theorem}[Perturbative Riccati-barrier regime]\label{thm:PerturbativeRiccatiBarrier}
Assume that the collar geometry, the MOTS stability data \((X,Q_\Sigma)\), and the logarithmic Green data \((\rho,\Upsilon)\) are \(C^{1,\alpha}\)-close to a model configuration for which
\begin{equation}\label{eq:ModelRiccatiGap37}
    \partial_\nu\rho_0+\rho_0^2+H_0\rho_0
    +2X_0\cdot\nabla_\Sigma\Upsilon_0
    +Q_0-\frac18 R_0
    -\frac18\mathcal E_0\mathbf 1
    \ge \kappa_0>0 .
\end{equation}
More precisely, suppose the \(C^{0,\alpha}\)-norm of the difference between the left side of~\eqref{eq:RiccatiBarrier37} and the model left side in~\eqref{eq:ModelRiccatiGap37} is at most \(\varepsilon\). If \(\varepsilon\le\kappa_0/2\), then the Green-trace calibration holds.
\end{theorem}

\begin{proof}
The assumed \(C^{0,\alpha}\) control gives the lower bound
\[
    \partial_\nu\rho+\rho^2+H_\Sigma\rho
    +2X\cdot\nabla_\Sigma\Upsilon
    +Q_\Sigma-\frac18 R_{\tg}^{\mathrm{reg}}\big|_\Sigma
    -\frac18\mathcal E_{\mathrm{collar}}\mathbf 1
    \ge \kappa_0-\varepsilon
    \ge \kappa_0/2 .
\]
Thus the Riccati barrier inequality holds. Proposition~\ref{prop:RiccatiBarrierCalibration} proves the claim.
\end{proof}

\begin{remark}[The next nonperturbative bridge]\label{rem:NextBridge37}
The remaining estimate is now first-order in the logarithmic normal flux. A viable route is to prove a Hopf--Riccati lower bound for \(\rho\) on the maximum-area MOTS, with the DEC and the MOTS stability potential controlling the zeroth-order terms in~\eqref{eq:RiccatiBarrier37}. This is more flexible than trying to force \(\tr_\Sigma k\ge0\): it only asks the ADM Green function to enter the horizon through the stability cone in the mass-relevant direction.
\end{remark}

\subsection{Dual Riccati calibration}
\label{subsec:DualRiccatiCalibration}

There is an even weaker target than pointwise flux-defect positivity. The conformal mass correction does not need \(L_\Sigma\Gamma_\Sigma\ge0\) everywhere; it only needs the pairing with the adjoint multiplier produced by the constrained maximum-area problem to have the correct sign.

\begin{definition}[Adjoint multiplier and weighted flux]\label{def:AdjointMultiplier37}
Let
\[
    h_\Sigma:=\tr_\Sigma k .
\]
For a strictly stable MOTS define the nonnegative adjoint multiplier
\[
    \zeta_\Sigma:=-L_\Sigma^{-*}h_\Sigma .
\]
Thus
\[
    L_\Sigma^*\zeta_\Sigma=-h_\Sigma,\qquad \zeta_\Sigma\ge0
\]
in the weak sense, where the inequality is exactly the adjoint-cone Euler condition of Theorem~\ref{thm:AdjointConeEuler}. We say that the ADM Green potential satisfies the \emph{dual Green-flux calibration} if
\begin{equation}\label{eq:DualFluxCalibration37}
    \int_\Sigma \zeta_\Sigma\,\mathfrak D_\Sigma(G_\infty)\,dA\ge0 .
\end{equation}
\end{definition}

\begin{theorem}[Dual calibration is enough for the mass sign]\label{thm:DualCalibrationMass}
Assume \(\Sigma\) is a strictly stable constrained maximum-area MOTS satisfying the adjoint-cone Euler condition. If the dual Green-flux calibration~\eqref{eq:DualFluxCalibration37} holds, then
\[
    \int_\Sigma(\tr_\Sigma k)\Gamma_\Sigma\,dA\le0,
\]
and hence the interface contribution to the ADM mass correction is nonpositive.
\end{theorem}

\begin{proof}
By definition, \(L_\Sigma^*\zeta_\Sigma=-\tr_\Sigma k\). Using the boundary identity \(L_\Sigma\Gamma_\Sigma=\mathfrak D_\Sigma(G_\infty)\), we compute
\[
    \int_\Sigma(\tr_\Sigma k)\Gamma_\Sigma\,dA
    =
    -\int_\Sigma (L_\Sigma^*\zeta_\Sigma)\Gamma_\Sigma\,dA
    =
    -\int_\Sigma \zeta_\Sigma L_\Sigma\Gamma_\Sigma\,dA
    =
    -\int_\Sigma \zeta_\Sigma\mathfrak D_\Sigma(G_\infty)\,dA .
\]
The assumed weighted flux inequality makes the final expression nonpositive. Lemma~\ref{lem:InterfaceADMGreen} then gives the desired ADM sign.
\end{proof}

\begin{corollary}[Penrose from dual calibration]\label{cor:PenroseFromDualCalibration}
In Theorem~\ref{thm:FullPenroseFromCalibration}, the Green-trace calibration hypothesis may be replaced by the dual Green-flux calibration~\eqref{eq:DualFluxCalibration37}.
\end{corollary}

\begin{proof}
The proof of Theorem~\ref{thm:FullPenroseFromCalibration} uses Green-trace calibration only through the sign of the interface mass contribution. Theorem~\ref{thm:DualCalibrationMass} supplies exactly that sign.
\end{proof}

Writing the flux defect in logarithmic variables gives the weighted Riccati form of the same target. Let
\begin{equation}\label{eq:RiccatiBracket37}
    \mathcal B_\Sigma(G_\infty)
    :=
    \partial_\nu\rho+\rho^2+H_\Sigma\rho
    +2X\cdot\nabla_\Sigma\Upsilon
    +Q_\Sigma-\frac18 R_{\tg}^{\mathrm{reg}}\big|_\Sigma
    -\frac18\mathcal E_{\mathrm{collar}}\mathbf 1 .
\end{equation}
Then \(\mathfrak D_\Sigma(G_\infty)=\Gamma_\Sigma\mathcal B_\Sigma(G_\infty)\), and the dual calibration becomes
\begin{equation}\label{eq:WeightedRiccati37}
    \int_\Sigma \zeta_\Sigma\Gamma_\Sigma
    \mathcal B_\Sigma(G_\infty)\,dA\ge0 .
\end{equation}

\begin{proposition}[Weighted Riccati barrier]\label{prop:WeightedRiccatiBarrier}
If the weighted Riccati inequality~\eqref{eq:WeightedRiccati37} holds, then the interface contribution to the ADM mass correction is nonpositive.
\end{proposition}

\begin{proof}
This is Theorem~\ref{thm:DualCalibrationMass} after substituting the logarithmic identity~\eqref{eq:LogFluxDefect37}.
\end{proof}

\begin{lemma}[Weighted logarithmic cancellation]\label{lem:WeightedLogCancellation}
Let
\[
    \mathcal R_\Sigma^{\mathrm{coll}}
    :=
    \frac18\left(\mathcal E_{\mathrm{collar}}^{\log}(\Phi)
    -\mathcal E_{\mathrm{collar}}\mathbf 1\right).
\]
Then the weighted Riccati pairing has the purely tangential form
\begin{equation}\label{eq:WeightedLogCancellation37}
    \int_\Sigma \zeta_\Sigma\Gamma_\Sigma
    \mathcal B_\Sigma(G_\infty)\,dA
    =
    \int_\Sigma \Gamma_\Sigma
    \left[
        \nabla_\Sigma\zeta_\Sigma\cdot\nabla_\Sigma\Upsilon
        +2\zeta_\Sigma X\cdot\nabla_\Sigma\Upsilon
        +\zeta_\Sigma\left(Q_\Sigma+\mathcal R_\Sigma^{\mathrm{coll}}\right)
    \right]dA .
\end{equation}
\end{lemma}

\begin{proof}
Use the logarithmic Riccati equation~\eqref{eq:GreenRiccati37} in the definition of \(\mathcal B_\Sigma\). The curvature terms \(R_{\tg}^{\mathrm{reg}}/8\) cancel, leaving
\[
    \mathcal B_\Sigma
    =
    -\Delta_\Sigma\Upsilon
    -|\nabla_\Sigma\Upsilon|^2
    +2X\cdot\nabla_\Sigma\Upsilon
    +Q_\Sigma
    +\mathcal R_\Sigma^{\mathrm{coll}} .
\]
Since \(\Sigma\) is closed,
\[
    -\int_\Sigma \zeta_\Sigma\Gamma_\Sigma\Delta_\Sigma\Upsilon\,dA
    =
    \int_\Sigma \nabla_\Sigma(\zeta_\Sigma\Gamma_\Sigma)\cdot\nabla_\Sigma\Upsilon\,dA .
\]
Using \(\nabla_\Sigma\Gamma_\Sigma=\Gamma_\Sigma\nabla_\Sigma\Upsilon\), the right side becomes
\[
    \int_\Sigma \Gamma_\Sigma\nabla_\Sigma\zeta_\Sigma\cdot\nabla_\Sigma\Upsilon\,dA
    +
    \int_\Sigma \zeta_\Sigma\Gamma_\Sigma|\nabla_\Sigma\Upsilon|^2\,dA .
\]
This cancels the term \(-\int_\Sigma\zeta_\Sigma\Gamma_\Sigma|\nabla_\Sigma\Upsilon|^2\,dA\), giving~\eqref{eq:WeightedLogCancellation37}.
\end{proof}

\begin{theorem}[Tangential drift-alignment criterion]\label{thm:TangentialDriftAlignment}
Assume \(\zeta_\Sigma\ge0\) is the adjoint multiplier from Definition~\ref{def:AdjointMultiplier37}. If
\begin{equation}\label{eq:TangentialAlignment37}
    \nabla_\Sigma\zeta_\Sigma\cdot\nabla_\Sigma\Upsilon
    +2\zeta_\Sigma X\cdot\nabla_\Sigma\Upsilon
    +\zeta_\Sigma\left(Q_\Sigma+\mathcal R_\Sigma^{\mathrm{coll}}\right)
    \ge0
\end{equation}
weakly on \(\Sigma\), then the weighted Riccati positivity conjecture holds. Consequently, after the compactness and standard smoothing/AMO hypotheses, the Penrose conclusion follows.
\end{theorem}

\begin{proof}
Multiplying~\eqref{eq:TangentialAlignment37} by the positive factor \(\Gamma_\Sigma\) and integrating gives a nonnegative right side in~\eqref{eq:WeightedLogCancellation37}. Therefore the weighted Riccati inequality~\eqref{eq:WeightedRiccati37} holds. Proposition~\ref{prop:WeightedRiccatiBarrier} gives the nonpositive interface mass contribution, and the completion is Corollary~\ref{cor:PenroseFromDualCalibration}.
\end{proof}

\begin{remark}[What remains after the cancellation]\label{rem:AfterCancellation37}
The second normal derivative and the quadratic tangential term \(|\nabla_\Sigma\Upsilon|^2\) have disappeared from the weighted bridge. The remaining target is a first-order tangential alignment between the Green trace \(\Upsilon\), the adjoint multiplier \(\zeta_\Sigma\), and the MOTS drift \(X\), with only the zeroth-order stability potential \(Q_\Sigma\) and the collar residual left to control. This is the next concrete estimate: prove~\eqref{eq:TangentialAlignment37}, or prove its integrated version, using the maximum-area Euler equation and the DEC-controlled structure of \(Q_\Sigma\).
\end{remark}

\begin{lemma}[Adjoint collapse of the tangential bridge]\label{lem:AdjointCollapse37}
In the notation above,
\begin{equation}\label{eq:AdjointCollapse37}
    \int_\Sigma \Gamma_\Sigma
    \left[
        \nabla_\Sigma\zeta_\Sigma\cdot\nabla_\Sigma\Upsilon
        +2\zeta_\Sigma X\cdot\nabla_\Sigma\Upsilon
        +\zeta_\Sigma\left(Q_\Sigma+\mathcal R_\Sigma^{\mathrm{coll}}\right)
    \right]dA
    =
    -\int_\Sigma (\tr_\Sigma k)\Gamma_\Sigma\,dA
    +
    \int_\Sigma \zeta_\Sigma\Gamma_\Sigma
        \mathcal R_\Sigma^{\mathrm{coll}}\,dA .
\end{equation}
\end{lemma}

\begin{proof}
Since \(\nabla_\Sigma\Gamma_\Sigma=\Gamma_\Sigma\nabla_\Sigma\Upsilon\), the first two terms on the left of~\eqref{eq:AdjointCollapse37} equal
\[
    \int_\Sigma
    \left(\nabla_\Sigma\zeta_\Sigma+2\zeta_\Sigma X\right)\cdot\nabla_\Sigma\Gamma_\Sigma\,dA .
\]
Integrating by parts on the closed surface \(\Sigma\) gives
\[
    -\int_\Sigma \Gamma_\Sigma
    \left(\Delta_\Sigma\zeta_\Sigma+2\div_\Sigma(\zeta_\Sigma X)\right)dA .
\]
Adding the zeroth-order term \(\int_\Sigma\zeta_\Sigma Q_\Sigma\Gamma_\Sigma\,dA\) produces
\[
    \int_\Sigma \Gamma_\Sigma
    \left(-\Delta_\Sigma\zeta_\Sigma-2\div_\Sigma(\zeta_\Sigma X)+Q_\Sigma\zeta_\Sigma\right)dA
    =
    \int_\Sigma \Gamma_\Sigma L_\Sigma^*\zeta_\Sigma\,dA .
\]
By the definition of the adjoint multiplier, \(L_\Sigma^*\zeta_\Sigma=-\tr_\Sigma k\). The remaining term is exactly the collar-residual contribution in~\eqref{eq:AdjointCollapse37}.
\end{proof}

\begin{proposition}[Residual-dominance closure]\label{prop:ResidualDominanceClosure}
Assume the adjoint multiplier \(\zeta_\Sigma\) is nonnegative and the collar residual satisfies
\begin{equation}\label{eq:ResidualDominance37}
    \int_\Sigma \zeta_\Sigma\Gamma_\Sigma
        \mathcal R_\Sigma^{\mathrm{coll}}\,dA
    \ge
    \int_\Sigma (\tr_\Sigma k)\Gamma_\Sigma\,dA .
\end{equation}
Then the weighted Riccati positivity conjecture holds. Consequently, under the compactness and standard smoothing/AMO hypotheses, the spacetime Penrose inequality follows.
\end{proposition}

\begin{proof}
The residual-dominance inequality~\eqref{eq:ResidualDominance37} makes the right side of~\eqref{eq:AdjointCollapse37} nonnegative. The weighted logarithmic cancellation therefore gives
\[
    \int_\Sigma \zeta_\Sigma\Gamma_\Sigma
    \mathcal B_\Sigma(G_\infty)\,dA\ge0 .
\]
This is weighted Riccati positivity. Theorem~\ref{thm:PenroseFromWeightedRiccati} completes the argument.
\end{proof}

\begin{remark}[Noncircular target]\label{rem:NonCircularTarget37}
Lemma~\ref{lem:AdjointCollapse37} is a useful diagnostic. If the exact collar model has \(\mathcal R_\Sigma^{\mathrm{coll}}=0\), then the weighted tangential criterion is equivalent to the desired Green-pairing sign and supplies no independent positivity by itself. A noncircular proof along this branch must therefore extract positive information from the collar residual, from a strengthened maximum-area Euler inequality, or from an independent comparison principle for the ADM Green trace.
\end{remark}

\begin{theorem}[Localized negative-defect absorption]\label{thm:LocalizedNegativeDefectAbsorption}
Let \(W_\Sigma:=\zeta_\Sigma\Gamma_\Sigma\). Suppose the Riccati bracket
\(\mathcal B_\Sigma=\mathcal B_\Sigma(G_\infty)\) satisfies
\begin{equation}\label{eq:NegativeDefectAbsorption37}
    \int_{\{\mathcal B_\Sigma<0\}} W_\Sigma|\mathcal B_\Sigma|\,dA
    \le
    \int_{\{\mathcal B_\Sigma>0\}} W_\Sigma\mathcal B_\Sigma\,dA .
\end{equation}
Then the dual Green-flux calibration holds, and therefore the Penrose conclusion follows once the compactness and standard smoothing/AMO hypotheses are available.
\end{theorem}

\begin{proof}
Splitting the weighted integral into its positive and negative parts gives
\[
    \int_\Sigma W_\Sigma\mathcal B_\Sigma\,dA
    =
    \int_{\{\mathcal B_\Sigma>0\}}W_\Sigma\mathcal B_\Sigma\,dA
    -
    \int_{\{\mathcal B_\Sigma<0\}}W_\Sigma|\mathcal B_\Sigma|\,dA .
\]
The absorption hypothesis~\eqref{eq:NegativeDefectAbsorption37} makes this integral nonnegative. Thus~\eqref{eq:WeightedRiccati37} holds, and Proposition~\ref{prop:WeightedRiccatiBarrier} gives the mass sign. The rest is the same completion argument as Corollary~\ref{cor:PenroseFromDualCalibration}.
\end{proof}

\begin{remark}[Why the dual target is strictly weaker]\label{rem:DualTargetWeaker37}
Pointwise flux-defect positivity requires \(\mathcal B_\Sigma\ge0\) everywhere. The dual Riccati calibration allows \(\mathcal B_\Sigma\) to be negative on regions where the adjoint multiplier \(\zeta_\Sigma\) gives little weight, provided the positive part dominates in the weighted balance~\eqref{eq:NegativeDefectAbsorption37}. This is the natural next estimate to pursue: the maximum-area variational problem already supplies the weight \(\zeta_\Sigma\), so the remaining geometric task is to prove that the ADM Green flux has nonnegative weighted Riccati average.
\end{remark}

\subsection{Nonperturbative target estimate}

The preceding theorems prove the route in genuine open regimes around strictly positive product cylinders and reduce the general case to either a pointwise maximum-principle estimate for the flux defect or the weaker weighted Riccati estimate attached to the adjoint multiplier.

\begin{conjecture}[Flux-defect positivity]\label{conj:FluxDefectPositivity}
Let \(\Sigma_{\max}\) be a strictly stable maximum-area MOTS obtained from the trapped-surface compactness program. Then the ADM Green potential satisfies
\[
    \mathfrak D_\Sigma(G_\infty)\ge0
\]
weakly on \(\Sigma_{\max}\).
\end{conjecture}

\begin{remark}[Why this is stronger than the previous roadmap]
Conjecture~\ref{conj:FluxDefectPositivity} is not a vague sign bridge. It is an explicit second-order boundary inequality for one distinguished positive solution \(G_\infty\). The numerical constants in \(\mathfrak D_\Sigma\) are tied to the conformal-Laplacian normalization fixed above; changing that normalization rescales the regular-curvature term but leaves the criterion unchanged after the corresponding definition is adjusted. The estimate can be attacked by Hopf boundary estimates, Riccati equations in the Jang collar, and the DEC-controlled zeroth-order terms in the MOTS stability operator.
\end{remark}

\begin{theorem}[Penrose from flux-defect positivity]\label{thm:PenroseFromFluxDefect}
Assume the compactness theorem for maximum-area trapped surfaces and the flux-defect positivity conjecture. Then the spacetime Penrose inequality holds for all closed future trapped surfaces covered by the compactness theorem, without a pointwise favorable-jump hypothesis.
\end{theorem}

\begin{proof}
Flux-defect positivity implies Green-trace calibration by Corollary~\ref{cor:FluxDefectCriterion}. The conclusion is then exactly Theorem~\ref{thm:FullPenroseFromCalibration}.
\end{proof}

\begin{conjecture}[Weighted Riccati positivity]\label{conj:WeightedRiccatiPositivity}
Let \(\Sigma_{\max}\) be a strictly stable maximum-area MOTS obtained from the trapped-surface compactness program. Then its adjoint multiplier \(\zeta_\Sigma=-L_\Sigma^{-*}(\tr_\Sigma k)\) and the ADM Green potential satisfy
\[
    \int_{\Sigma_{\max}}\zeta_\Sigma\Gamma_\Sigma
    \mathcal B_\Sigma(G_\infty)\,dA\ge0 .
\]
\end{conjecture}

\begin{theorem}[Penrose from weighted Riccati positivity]\label{thm:PenroseFromWeightedRiccati}
Assume the compactness theorem for maximum-area trapped surfaces and the weighted Riccati positivity conjecture. Then the spacetime Penrose inequality holds for all closed future trapped surfaces covered by the compactness theorem, without a pointwise favorable-jump hypothesis.
\end{theorem}

\begin{proof}
Weighted Riccati positivity is exactly the dual Green-flux calibration in logarithmic form. The conclusion follows from Corollary~\ref{cor:PenroseFromDualCalibration}.
\end{proof}
